\chapter{My Journey at UFV}
\setcounter{section}{0}

\section{The PFG as the Culmination of My University Journey}

When I first arrived at the university, I did not have much experience with programming, but I had always felt curiosity about computing from a practical point of view. Before starting the degree, I had spent time with small hardware hobbies, such as disassembling mobile phones and computers to understand how they worked inside. What attracted me most was the possibility of creating things that could make everyday life easier. I did not begin university with very clear expectations, but during the first year, especially through the mentoring sessions, I discovered how much I enjoyed Computer Science. I remember sharing with my mentor the projects I had in mind and, that summer, turning many of those ideas into real projects.

My university journey was not a straight line. Some subjects motivated me more than others, but courses such as Object-Oriented Programming, Computer Architecture, Data Structures and Algorithms gave me a strong technical foundation and helped me see computing as something creative, rigorous, and useful. The biggest change came during my third year, when I went on Erasmus. Although passing some subjects was more difficult, I also had more time to investigate on my own. This coincided with the rise of artificial intelligence, and I started learning more about AI, Linux, and the tools that allow a computer system to be understood and controlled more deeply.

In the second semester of third year, subjects such as Artificial Intelligence II and Data Mining consolidated that interest. My tutor, Mois{\'e}s Mart{\'i}nez, was very important throughout this path, as he taught me during almost every year of the degree and later guided me in areas directly connected to this project. My classmates also pushed me to believe more in my own work, learn more, and try to achieve better results. Looking back, I changed from someone who only suspected that computing might be what he liked into someone who genuinely enjoys building tools with a computer, both to make his own life easier and to help others.

Now I see myself with much more knowledge than when I started, but also with a greater awareness of how much I still do not know. At the beginning of the degree, I felt that I was learning a lot and that I knew a lot. Now I know much more, but I also see the scale of the field and the quality of the work being done inside and outside the university. That has made me more humble, more committed, and more constant.

The motivation for this PFG came from a question I had been thinking about for some time: how could I automate my life and make everyday tasks faster and easier? With the rise of AI, I began to think that artificial intelligence should not remain only as a chatbot that answers questions. I was interested in an assistant integrated into a user's digital life, capable of receiving an objective and reasoning through the steps needed to complete it. I imagined a system that could help with tasks such as booking a flight, searching for information, configuring something on the computer, writing a document, opening files, organizing tasks, or teaching the user how to do something new.

At first, I considered creating a Linux distribution based on AI. It was an ambitious idea, but too costly and difficult to turn into a visible result within the scope of the PFG. That idea evolved into Autopilot: a system that could be installed and used to perform desktop tasks through voice. The project relates closely to who I am because I have always liked creating tools that make life easier. It connects my curiosity for AI, Linux, automation, and personal projects with a practical goal.

The approach was shaped by reliability and privacy. An AI that automates tasks can be very useful, but if the task is critical, hallucinations cannot be accepted. For that reason, I chose a neuro-symbolic approach, combining AI capabilities with a more structured planning layer. A personal computer also contains sensitive information, so Autopilot was designed as a local, on-device system. Since local computation is limited, the architecture followed a *divide and conquer* strategy: instead of relying on one large monolithic system, the problem was divided into smaller agents and modules with specific responsibilities. I also considered real usefulness, accessibility, feasibility, computational cost, modularity, and the possibility of evaluating the system against other alternatives.

\section{Connection with My Professional Future}

Completing this project has helped me clarify the type of professional path I want to follow. Autopilot shows that I am interested in technology with a purpose, not only in technology because it is new. The project is aligned with values that are important to me, such as *open source*, data privacy, and the creation of tools that can be useful for people who need them. I believe that users are increasingly exposed in the digital world, and that software should respect their personal information. In this case, the project could also help people with reduced mobility, allowing them to execute tasks on a computer using only their voice.

This PFG has also raised questions about what kind of engineer I want to become. Working on AI agents made me reflect on the balance between autonomy and control, and on the responsibility involved when a system can act on behalf of a user. I discovered that I am especially interested in building AI systems that are useful, private, modular, and understandable. I also enjoy the space between research and engineering: learning new ideas, testing them, and turning them into something that can be measured and used.

After these years at university, I can see myself working in areas related to artificial intelligence, software engineering, automation, accessibility, or applied research. The project has opened the perspective of building systems that do not simply answer, but act responsibly and help people interact with technology in a more natural way. Moving forward, I want to keep improving my technical skills and participate in projects where technology has a clear practical or social impact.

For that reason, this PFG is not only the end of my university journey. It is also a bridge toward the kind of work I want to pursue: technically demanding, useful for real users, respectful of privacy, and guided by the question of how technology can make people's lives easier in a responsible way.
